\documentclass[11pt]{article}
\usepackage[utf8]{inputenc}
\usepackage[english]{babel}
\usepackage{geometry}
\usepackage{amsmath}
\usepackage{booktabs}
\usepackage{xcolor}
\usepackage{mdframed}

\geometry{a4paper, margin=1in}

\title{Technical Specifications: Ettus USRP N210}
\author{Radiofrequency Engineering}
\date{\today}

\newmdenv[
  backgroundcolor=red!8,
  linecolor=red!60,
  linewidth=1pt,
  roundcorner=4pt,
  innertopmargin=6pt,
  innerbottommargin=6pt
]{warningbox}

\newmdenv[
  backgroundcolor=yellow!10,
  linecolor=yellow!60!black,
  linewidth=1pt,
  roundcorner=4pt,
  innertopmargin=6pt,
  innerbottommargin=6pt
]{pendingbox}

\begin{document}

\maketitle
\tableofcontents
\newpage

% =========================================================
\section{Gains}
% =========================================================

Controlled through the USRP Hardware Driver (\texttt{UHD}). They depend directly on the installed daughterboard (e.g., UBX, CBX, SBX):

\begin{itemize}
    \item \textbf{Global RX Gain:} \texttt{usrp->set\_rx\_gain()} - The range and steps depend on the board. On a typical UBX-40, it goes from 0 to 31.5 dB.
    \item \textbf{Global TX Gain:} \texttt{usrp->set\_tx\_gain()} - On a UBX-40, typically from 0 to 31.5 dB.
    \item \textbf{Specific Elements (PGA/VGA):} \texttt{usrp->set\_rx\_gain(val, "PGA0")} - Allows adjustment of individual amplification stages if the daughterboard supports it.
\end{itemize}

\textbf{Note:} To maximize dynamic range without saturating the ADC (14-bit), the clipping indicator in UHD should be monitored (\texttt{O} for overruns, although analog saturation requires observing the baseband spectrum).

% =========================================================
\section{Hardware Filters}
% =========================================================

\begin{itemize}
    \item \textbf{Analog Filters (Daughterboard):} Low-pass filters (LPF) or band-pass filters configured automatically by UHD according to the frequency and requested analog bandwidth (\texttt{set\_rx\_bandwidth()}).
    \item \textbf{Digital Processing (Spartan 3A-DSP 3400 FPGA):}\\
    Implements DDC (Digital Down Converter) in RX and DUC (Digital Up Converter) in TX with 25 MHz resolution. Handles decimation and interpolation from the motherboard's fixed 100 MHz master clock, using CIC filters followed by compensating FIR filters (half-band) to prevent aliasing.
\end{itemize}

% =========================================================
\section{Frequency Ranges}
% =========================================================

\begin{itemize}
    \item \textbf{Base Range (Motherboard):} DC to 100 MHz. Handles baseband or pure IF signals directly if LFRX/LFTX boards are used.
    \item \textbf{RF Range (Depends on Daughterboard):}
    \begin{itemize}
        \item \textbf{UBX-40:} 10 MHz to 6 GHz.
        \item \textbf{CBX:} 1.2 GHz to 6 GHz.
        \item \textbf{SBX:} 400 MHz to 4.4 GHz.
    \end{itemize}
    \item The motherboard supports fine FPGA-side frequency adjustment (\textit{DSP tuning}) using CORDICs. Frequency changes can be nearly instantaneous because they avoid the RF synthesizer PLL settling delay.
\end{itemize}

% =========================================================
\section{Sample Rate and Resolution}
% =========================================================

Limited by the Gigabit Ethernet interface and the motherboard converters: 100 MS/s ADC (14-bit) and 400 MS/s DAC (16-bit). The board clock is \textbf{fixed at 100 MHz} - unlike other USRPs, it is not user-configurable.

\begin{table}[h]
    \centering
    \small
    \begin{tabular}{@{}p{0.24\textwidth}p{0.20\textwidth}p{0.22\textwidth}p{0.25\textwidth}@{}}
        \toprule
        \textbf{Format} & \textbf{I/Q Resolution} & \textbf{Max. Sample Rate} & \textbf{Bandwidth} \\
        \midrule
        \texttt{sc16} (Standard) & 16-bit & 25 MS/s & 25 MHz observable \\
        \texttt{sc8} (High speed) & 8-bit & 50 MS/s & 50 MHz observable \\
        \texttt{sc16} (MIMO) & 16-bit & $\leq$ 12.5 MS/s/ch & 2 channels; practical Ethernet-bandwidth limit \\
        \bottomrule
    \end{tabular}
\end{table}

\textbf{Note:} Using the \texttt{sc8} wire format (\texttt{otw\_format}) allows the bandwidth to be doubled while losing dynamic resolution, expanding the signal again to 16- or 32-bit floating point on the host PC.

% =========================================================
\section{IQ \& DC Corrections}
% =========================================================

\begin{itemize}
    \item \textbf{DC Offset \& LO Leakage:} UHD performs dynamic DC tracking and cancellation by default in the FPGA.
    \item \textbf{I/Q Imbalance:}
    \begin{pendingbox}
    Although mitigated by hardware, I/Q imbalance in the RF front end requires calibration for critical applications. Use the UHD calibration utilities to generate and save compensation parameters in the EEPROM:

    \noindent\texttt{uhd\_cal\_rx\_iq\_balance}\\
    \texttt{uhd\_cal\_tx\_iq\_balance}
    \end{pendingbox}
\end{itemize}

% =========================================================
\section{Connectors and Inputs}
% =========================================================

\begin{table}[h]
    \centering
    \small
    \begin{tabular}{p{0.18\textwidth}p{0.26\textwidth}p{0.44\textwidth}}
        \toprule
        \textbf{Connector} & \textbf{Function} & \textbf{Specification} \\
        \midrule
        Gigabit Ethernet & Data & RJ45, 1000BASE-T. MTU 1500 recommended (9000 for Jumbo). Default IP: 192.168.10.2. \\
        Power & DC current & 6 V DC (included supply: 3 A max.). Typical consumption: 1.3 A at idle, 2.3 A with WBX. \\
        REF IN (SMA) & Reference clock & 10 MHz. Recommended level: 0 to +15 dBm. Square wave preferred; sine wave acceptable. \\
        PPS IN (SMA) & Trigger / synchronization & Pulse per second (1 PPS). Accepts TTL 0--3.3 V and 0--5 V. \\
        MIMO Cable & N210 interconnection & Shares clock and data between master and slave. Forms a coherent $2\times2$ system. \\
        RF (SMA) & TX/RX (on Daughterboard) & 50 $\Omega$. TX/RX and RX2 ports depending on installed DB. \\
        \bottomrule
    \end{tabular}
\end{table}

% =========================================================
\section{Safe Operating Limits}
% =========================================================

\begin{warningbox}
\textbf{Critical warnings - the N210 and its daughterboards are sensitive}
\end{warningbox}

\vspace{4pt}
\begin{itemize}
    \item \textbf{Maximum RX power:} Do not exceed \textbf{+15 dBm} on the RF port of any daughterboard (less on some boards such as the WBX). Permanent LNA damage is guaranteed above this threshold. For optimal linearity, keep below -15 dBm.
    \item \textbf{TX$\to$RX loopback:} Never connect TX directly to RX without an attenuator of at least 30 dB. Typical output power with a WBX is around \textbf{+15 dBm} - not the +20 dBm that some unofficial source may indicate. The value depends on the daughterboard and frequency.
    \item \textbf{Thermal Load:} The fan must be operational. The chassis dissipates heat from the ADC, DAC, and FPGA. Do not block the side vents.
    \item \textbf{Synchronization Voltages:} Do not inject more than +15 dBm into REF IN. PPS IN accepts a maximum of 5 V TTL.
\end{itemize}

% =========================================================
\section{Clock and Synchronization}
% =========================================================

\begin{itemize}
    \item \textbf{Internal Clock (TCXO):} \textbf{2.5 ppm} reference oscillator, active by default. The motherboard master clock is fixed at 100 MHz and is not user-configurable.
    \item \textbf{GPSDO (Optional):} If the GPS-disciplined module is installed, it improves accuracy to \textbf{0.01 ppm}. UHD detects it automatically; it requires writing a value to the EEPROM to confirm it: \texttt{usrp\_burn\_mb\_eeprom --values="gpsdo=internal"}.
    \item \textbf{External Clock:} 10 MHz signal through the front-panel SMA (REF IN). Configure in UHD with \texttt{usrp->set\_clock\_source("external")}. A square wave provides better phase noise; a sine wave is acceptable.
    \item \textbf{Time Synchronization (1 PPS):} Fundamental for aligning captures from multiple N210 units in the time domain. Configure with \texttt{usrp->set\_time\_source("external")} together with \texttt{set\_time\_next\_pps()}.
    \item \textbf{MIMO Cable:} Allows one USRP (slave) to share another USRP's (master) master clock and Ethernet network. On the slave device: \texttt{usrp->set\_clock\_source("mimo")} and \texttt{usrp->set\_time\_source("mimo")}. The external clock should only be supplied to the master.
\end{itemize}

% =========================================================
\section{Indicator LEDs}
% =========================================================

Located on the front panel. Vital for diagnosing connectivity and PLL status without needing software.

\begin{table}[h]
    \centering
    \begin{tabular}{p{0.08\textwidth}p{0.28\textwidth}p{0.52\textwidth}}
        \toprule
        \textbf{LED} & \textbf{Function} & \textbf{State and interpretation} \\
        \midrule
        \textbf{A} & TX activity & On/blinking during active transmission. \\
        \addlinespace
        \textbf{B} & RX activity & On/blinking during active reception. \\
        \addlinespace
        \textbf{C} & Firmware loaded & \textbf{On:} active firmware and programmed FPGA. \textbf{Off:} firmware not loaded - indicates a boot problem or incorrect image. \\
        \addlinespace
        \textbf{D} & Ethernet link & On or blinking when the Gigabit link is active. \\
        \addlinespace
        \textbf{E} & PLL Lock (Ref) & \textbf{Steady:} PLL locked to the master clock (internal or external). \textbf{Fast blinking:} PLL not locked - REF IN signal absent, weak, or out of frequency. Brief blinking at boot or when stopping a session is normal. \\
        \addlinespace
        \textbf{F} & PPS present & Blinks once per second if it receives a valid PPS signal. \\
        \bottomrule
    \end{tabular}
\end{table}

\textbf{Normal startup state (without an active session):} LEDs C, D, and E should be on. A and B off. F blinks only if a PPS signal is connected.

\textbf{Known pitfalls:}
\begin{itemize}
    \item If LED C is off after power-on, the firmware is not loaded. Use \texttt{uhd\_image\_loader} to reprogram.
    \item If LED E blinks rapidly and persistently during operation, the PLL has lost lock - check the REF IN signal.
    \item Do not load USRP2 images on the N210: they are incompatible and can leave the device inoperable.
\end{itemize}

% =========================================================
\section{Expansion Interfaces (GPIO and Control)}
% =========================================================

\subsection*{Daughterboard GPIO}
Certain boards (such as BasicRX/BasicTX or LFRX/LFTX) expose GPIO pins routable to the FPGA.
\begin{itemize}
    \item \textbf{UHD Control:} \texttt{set\_gpio\_attr()} in the UHD API.
    \item \textbf{Typical use:} Control of external antenna relays (T/R switches), digital attenuators, or power-amplifier (PA) switching. 3.3 V CMOS logic levels.
\end{itemize}

\subsection*{JTAG Port}
Standard internal Xilinx connector. Allows a bitstream to be dumped directly to the Spartan 3A-DSP FPGA while bypassing flash memory. Used exclusively for custom IP development with Xilinx ISE (tool compatible with Spartan 3).

\textbf{Note:} The N210 does not natively support RFNoC - that capability is exclusive to the X300/X310 series. To customize FPGA processing, modify and compile UHD's FPGA tree with Xilinx tools compatible with the Spartan 3 architecture.

\end{document}
